\section{Phép cắt gradient trong bản gốc}
\label{sec:ch04-cat-gradient}

\index{lan truyền ngược cắt cụt!ba phép thế}%
Tóm tắt của bài mở câu thứ ba bằng \enquote{Truncating the gradient where this
does not do harm}~\autocite{hochreiter1997lstm}. Phép cắt nằm ngay ở đó, trong
tóm tắt, câu liền sau câu giới thiệu cái tên LSTM. Thuật toán 1997 không tính
gradient; nó tính gradient với ba đạo hàm bị thay bằng không, và ba đạo hàm ấy
được liệt kê tên trong phụ lục A.1:

\begin{equation}
  \jac{\vect{a}^{i}_t}{\vect{h}_{t-1}} \to \mat{0}, \qquad
  \jac{\vect{a}^{o}_t}{\vect{h}_{t-1}} \to \mat{0}, \qquad
  \jac{\vect{a}^{c}_t}{\vect{h}_{t-1}} \to \mat{0}.
  \label{eq:ch04-ba-phep-the}
\end{equation}

Đọc kỹ ba dòng này thì thấy chúng cắt đúng một thứ: đường đi ngược \emph{qua}
\(\vect{h}_{t-1}\) về quá khứ. Bài nói rõ phần còn lại được giữ: lỗi tới
net input của \tn{ô nhớ}{memory cell} và của cổng không lan ngược xa hơn trong
thời gian, dù vẫn đổi được các trọng số vào (\enquote{do not get propagated
back further in time (although they \emph{do} serve to change the incoming
weights)})~\autocite{hochreiter1997lstm}.

Khác biệt ấy quan trọng, và trong code nó là một chữ:

\begin{minted}{python}
for t in range(inputs.shape[0]):
    h_prev = states[-1].h.detach() if truncate else states[-1].h
    states.append(self.step(states[-1].c, h_prev, inputs[t]))
\end{minted}

\code{detach} trên \(\vect{h}_{t-1}\) làm đúng việc bài mô tả: tích với
\(\vect{h}_{t-1}\) vẫn được lấy nên \(\mat{W}_{hi}\), \(\mat{W}_{ho}\) và
\(\mat{W}_{hc}\) vẫn nhận gradient của mình, nhưng không có gì chảy ngược tiếp
qua \(\vect{h}_{t-1}\). Chỗ \emph{không} được detach là \code{states[-1].c}.
Đường đó là CEC, và cắt nó là bỏ đi thứ duy nhất bài xây.

\subsection*{Phép cắt đắt bao nhiêu}

\index{lan truyền ngược cắt cụt!chi phí đo được}%
Bài khẳng định phép cắt không gây hại, và khẳng định ấy có chống lưng: họ nói
đã chạy vài thí nghiệm với LSTM không cắt và không thấy khác biệt đáng kể, với
lý do là ngoài CEC thì dòng lỗi vốn tắt rất nhanh. Đó là một khẳng định về
\emph{kết quả huấn luyện}. Tôi muốn một con số cho một thứ hẹp hơn và kiểm được
rẻ hơn: hai vector gradient ấy lệch nhau bao nhiêu về hướng.

\begin{measured}{góc giữa gradient cắt cụt và gradient đầy đủ}
16 đơn vị ẩn, trọng số trong \([-0.1, 0.1]\), loss ở cuối chuỗi, cosine lấy
trên cả chín tensor trọng số ghép lại. 20 lần rút mỗi độ dài, vì một lần rút
cho một con số chứ không cho một xu hướng:

\begin{minted}{text}
    T  mean cosine  min cosine   mean norm ratio
   10  0.991283     0.983478     0.985044
   25  0.946810     0.882489     0.941741
   50  0.804282     0.657828     0.740773
  100  0.451933     0.202544     0.567376
\end{minted}

Ở chuỗi 10 bước hai gradient gần như trùng nhau. Ở 100 bước cosine trung bình
còn 0.451933, và lần rút tệ nhất còn 0.202544.

Đo bằng \code{experiments/ch04\_truncation.py} tại tag \code{ch04}.
\end{measured}

Con số này không mâu thuẫn với bài, và chỗ nó không mâu thuẫn đáng nói rõ, vì
đọc lướt thì nó trông như một cái bẫy vừa sập. Bài nói về kết quả huấn luyện
sau khi chạy xong; tôi đo góc giữa hai vector tại khởi tạo. Hai phát biểu về
hai đại lượng khác nhau, và cả hai có thể cùng đúng.

Chỗ thú vị là ghép nó với số đo ở mục~\ref{sec:ch04-dao-ham}. Thành phần mà
phép cắt vứt đi chính là thành phần phình theo khoảng cách, thành phần đưa
\(\norm{\jac{\vect{c}_t}{\vect{c}_k}}\) lên 128.34 ở khoảng cách 100. Hai bảng
đo hai thứ nhưng chúng khớp vào nhau: gradient đầy đủ ở chuỗi dài lệch khỏi
gradient cắt cụt \emph{vì} nó bị thành phần rò rỉ ấy chi phối. Vậy nên
\enquote{cắt đi mà không hại} có một cách đọc mạnh hơn cách bài viết ra. Cái bị
cắt không phải thông tin bị hy sinh cho tốc độ. Nó là đường duy nhất trong ô
nhớ vẫn còn hành xử như một RNN thường, và cắt nó là chỗ lời hứa
\enquote{không đổi} chuyển từ gần đúng sang đúng. Con số 128.34 và bảng cosine
ở trên là hai cách nhìn cùng một thành phần.

Tôi không đọc chỗ này thành \enquote{bài trình bày sai lý do của chính mình}.
Bài nói phép cắt hiệu quả và không hại, và cả hai vế đều đúng. Chỉ là phần
\enquote{không hại} hóa ra không phải một sự nhân nhượng, và
phương trình~\eqref{eq:ch04-jac-cat} của mục trước không tồn tại nếu thiếu nó.

\subsection*{Cái phép cắt thật sự lấy đi}

\index{lan truyền ngược cắt cụt!giới hạn}%
Có một cái giá, và bài tự nêu ở mục 6 chứ không giấu. Bản cắt cụt không giải
được những bài kiểu \enquote{strongly delayed XOR}, nơi phải lấy XOR của hai
input cách xa nhau. Lý do không phải là dòng lỗi: lưu \emph{một} trong hai
input không làm giảm sai số kỳ vọng chút nào, nên bài toán không phân rã được
thành một mục tiêu con dễ hơn để bám vào. Mạng không có bậc thang nào để leo.

Bài nói thẳng rằng về lý thuyết dùng gradient đầy đủ có thể vòng qua chỗ này,
và nói luôn là họ không khuyến nghị: nó tốn hơn, dòng lỗi không đổi qua CEC chỉ
chứng minh được cho bản cắt cụt, và mấy thí nghiệm họ đã chạy với bản không cắt
không cho khác biệt đáng kể.
